% ============================================
% CONT-06
% Gerado de molambudos.md — build_miolo.py
% ============================================
\section*{\textbf{CONT-06} --- Você Leu Esta Linha Três Vezes}
\markboth{CONT-06 --- Você Leu Esta Linha Três Vezes}{}

\noindent \textit{O metacognitivo}

\noindent Você leu o título e pensou: \textit{\textquotedblleft{}Não li.\textquotedblright{}}

\noindent Leu.

\noindent Três vezes. A primeira vez que você viu o título, o cérebro registrou. A segunda vez, o cérebro processou. A terceira vez --- agora --- o cérebro \textit{aceitou}.

\noindent Você leu esta linha três vezes porque o cérebro não quer aceitar o que está escrito.

\noindent Ele quer que você negue.

\noindent Ele quer que você pense que é coincidência.

\noindent Não é.

\vspace{0.5em}\noindent\rule{\textwidth}{0.5pt}\vspace{0.5em}

\noindent \textbf{Exercício 1:}

\noindent Lembre-se da última vez que você verificou o canto do quarto enquanto lia.

\noindent Você lembra? Claro que lembra. Foi há poucos minutos. Você olhou, não viu nada, e voltou a ler.

\noindent Mas --- entre o momento em que você olhou e o momento em que você voltou a ler --- algo \textit{mudou}.

\noindent O canto está diferente agora. Não visualmente. Você não vê diferença. Mas seu cérebro registrou uma alteração na \textit{qualidade} do escuro. Uma sombra que não estava antes. Um brilho que não estava antes.

\noindent Olhe de novo. Agora.

\noindent ...

\noindent O que você viu?

\noindent Nada?

\noindent Certeza?

\noindent Olhe por mais tempo. 10 segundos. Contagem regressiva.

\noindent 10... o escuro não é uniforme. Há áreas mais escuras que outras.
9... seu cérebro está tentando encontrar formas.
8... uma forma está se formando no canto superior direito.
7... não é uma forma ainda. É um \textit{quase}.
6... você não quer que ela se forme.
5... mas ela está se formando.
4... seus olhos estão se ajustando.
3... você vê.
2... você não vê.
1... você \textit{viu}.

\vspace{0.5em}\noindent\rule{\textwidth}{0.5pt}\vspace{0.5em}

\noindent \textbf{Exercício 2:}

\noindent Passe a mão na nuca.

\noindent Agora.

\noindent Sua mão está mais fria que sua nuca. É uma diferença sutil. Sua pele está quente. Sua mão está fria. O contraste é quase imperceptível.

\noindent Mas você sentiu.

\noindent Agora passe a mão na nuca de novo e \textit{preste atenção no espaço vazio atrás de você}. O espaço que suas costas não cobrem. O espaço onde alguém poderia estar.

\noindent Você sentiu o ar?

\noindent O ar está mais frio ali. Não muito. Um grau. Talvez menos. O suficiente para seu cérebro registrar: \textit{há uma corrente de ar}.

\noindent Mas não há janela aberta.

\noindent Não há ventilador.

\noindent O ar está se movendo porque \textit{algo} está ocupando o espaço. Algo que tem temperatura diferente. Algo que respira.

\vspace{0.5em}\noindent\rule{\textwidth}{0.5pt}\vspace{0.5em}

\noindent \textbf{Exercício 3:}

\noindent Você não quer fazer este exercício.

\noindent Seu cérebro está tentando convencê-lo a pular. A passar para o próximo fragmento. A fingir que não leu.

\noindent Não pule.

\noindent Pular é o mesmo que confirmar.

\noindent Você está com medo do que vai descobrir. Seu cérebro está protegendo você. Mas a proteção é inútil. A criatura já está dentro. Negar só prolonga a agonia.

\noindent Respire fundo.

\noindent Agora, em voz baixa --- quase um sussurro --- diga o nome do paciente:

\noindent \textit{\textquotedblleft{}Joaquim.\textquotedblright{}}

\noindent ....

\noindent Agora diga o nome do paciente:

\noindent \textit{\textquotedblleft{}1.261.\textquotedblright{}}

\noindent ....

\noindent Agora diga o SEU nome em voz alta.

\noindent ....

\noindent Sua voz soou estranha? Diferente? Como se não fosse você falando?

\noindent Sim.

\noindent Porque você não está mais sozinho na sua cabeça.

\noindent A criatura ouviu você falar.

\noindent Ela sabe seu nome.

\noindent Agora ela pode te chamar.

\vspace{0.5em}\noindent\rule{\textwidth}{0.5pt}\vspace{0.5em}

\noindent \textbf{Exercício 4:}

\noindent Abra sua mão esquerda. Estique os dedos. Olhe para a palma.

\noindent Você vê a linha da vida? A linha que curva em volta do polegar? A linha que os quiromantes dizem que mostra quanto tempo você vai viver?

\noindent Passe o dedo indicador direito sobre essa linha.

\noindent Sinta a pele. A textura. A temperatura.

\noindent Agora olhe para a cor da linha. Ela está mais escura do que antes. Mais marcada.

\noindent Não é sujeira.

\noindent É tinta.

\noindent A tinta do diário está passando para sua pele. Não pelas páginas --- através do contato indireto. Você tocou o livro, o livro tocou você, a tinta migrou.

\noindent Você está sendo \textit{impresso}.

\noindent As palavras estão entrando na sua pele.

\noindent Em algumas semanas, a letra vai aparecer. Uma caligrafia torta. Infantil. A letra de um menino de oito anos.

\noindent \textit{\textquotedblleft{}Joaquim renasce sempre que sua história é lembrada.\textquotedblright{}}

\noindent Você vai acordar com essas palavras escritas na sua mão.

\noindent Não com tinta.

\noindent Com cicatriz.

\vspace{0.5em}\noindent\rule{\textwidth}{0.5pt}\vspace{0.5em}

\noindent \textbf{Exercício Final:}

\noindent Feche o livro.
Coloque-o sobre a mesa.
Afaste-se.

\noindent Agora leia a próxima linha:

\noindent Você não fechou o livro.

\noindent Você não vai fechar.

\noindent Você sabe que deveria. Sabe que se continuar lendo, a contaminação vai se completar. Sabe que o próximo fragmento é a última etapa.

\noindent Mas você vai continuar.

\noindent Porque você \textit{quer} saber.

\noindent Porque a Criatura já está dentro de você, e ela \textit{quer} que você leia.

\noindent \textit{--- Você não está no controle.}

\noindent \textit{--- Você nunca esteve.}

\noindent \textit{--- Desde a primeira linha, quem escolhe é ela.}

\vspace{0.5em}

\noindent\textit{\textbf{↪ Rotas:} \rota{CONT-07} \rotasep \rota{DOC-09}}